%% kapa: exam=1 round=29 year=2018 period=1 form=A subject=civil_law subject_ko=민법 seq=33 number=33
%% source: Q-Net 감정평가사 기출문제 (cst00309) artlSeq=5202808 file=extracted/1008/5202808/2018년 제29회 감정평가사 1차시험 기출문제/2018년 제29회 감정평가사 1차시험 기출문제/1교시 A형.pdf
%% conversion: kordoc 4.24.51 -> tools/gen1.py (auto); review=proofread-sonnet; flags=none
\begin{question}{33}
점유보호청구권에 관한 설명으로 옳지 않은 것은? (다툼이 있으면 판례에 따름)
\begin{choices}
  \item 점유물방해제거청구권의 행사기간은 출소기간이다.
  \item 점유보조자에게는 점유물방해제거청구권이 인정되지 않는다.
  \item 직접점유자가 임의로 점유를 타인에게 이전한 경우, 그 점유이전이 간접점유자의 의사에 반하더라도 간접점유자의 점유가 침탈된 경우에 해당하지 않는다.
  \item 점유자가 점유의 침탈을 당한 경우, 침탈자의 특별승계인이 악의인 때에도 그 특별승계인에게 점유물반환청구권을 행사할 수 없다.
  \item 공사로 인하여 점유의 방해를 받은 경우, 공사 착수 후 1년을 경과하거나 그 공사가 완성된 때에는 방해의 제거를 청구하지 못한다.
\end{choices}
\end{question}
\answer{4}
